\section{Coding and operational preliminaries}
\label{sec:core}
Let $H\in\Ftwo^{m\times n}$, with variable set $\Var=\{1,\ldots,n\}$,
check set $\Chk=\{1,\ldots,m\}$, and Tanner edges
$\Edges=\{(j,i):H_{ji}=1\}$. Write $\Nb(j)=\{i:H_{ji}=1\}$ and
$C=\{x\in\Ftwo^n:Hx=0\}$. Vectors are columns throughout.
The matrix need not have full row rank, and empty rows, empty columns,
$m=0$ and $n=0$ are allowed. Sparsity is needed for the implementation
bounds, not for the algebraic theorems. LDPC decoding and its standard
iterative and algebraic ingredients are classical
\cite{Gallager1962,RichardsonUrbanke2008,LubyMitzenmacherShokrollahiSpielman2001,BurshteinMiller2004ML}.

A received word is $y\in\{0,1,\starE\}^n$. A partial assignment
$K:\Var\rightharpoonup\Ftwo$ has unknown set
$U=\Var\setminus\operatorname{dom}K$. A full word $x$ extends $K$ if
$x_i=K(i)$ whenever $K(i)$ is defined. Initially $K_0(i)=y_i$ at each
nonerased coordinate. For an assignment $K$, define
\begin{equation}\label{eq:operational-state}
 s_K(j)=\sum_{i\in\Nb(j)\cap\operatorname{dom}K}K(i),
 \qquad d_K(j)=|\Nb(j)\cap U|.
\end{equation}
All sums and matrix products are in $\Ftwo$, unless explicitly stated
otherwise. Thus addition is XOR. A coherent operational state
$\sigma=(K,U,s,d)$ has $s=s_K$ and $d=d_K$.

\begin{defi}[Operational peeling and trace]\label{def:peeling}
A triple $a=(j,i,b)$ is valid at $\sigma$ if
$\Nb(j)\cap U=\{i\}$ and $b=s(j)$. Its successor extends $K$ by
$K(i)=b$, removes $i$ from $U$, and, for every $\ell$ adjacent to $i$,
sets $s(\ell)\leftarrow s(\ell)+b$ and
$d(\ell)\leftarrow d(\ell)-1$. Other entries remain unchanged.
A trace $T=(a_1,\ldots,a_t)$ is valid from $y$ when every triple is valid
at the preceding state, starting at $K_0$. Traces may be empty or stop
before peeling is maximal. Every valid trace has $t\leq n$ because it
resolves a different unknown coordinate at every step.
\end{defi}

\begin{lem}[State meaning and preservation of completions]\label{lem:core-invariant}
Every state reached by a valid trace satisfies
\eqref{eq:operational-state}. Furthermore, every $x\in C$ agreeing with
the visible coordinates of $y$ extends its reached assignment $K$.
\end{lem}
\begin{proof}
Both state identities hold initially by definition. In a valid step,
$i$ was unknown, so adding it to $\operatorname{dom}K$ adds exactly $b$
to the known sum at each adjacent check and removes exactly one term
from the corresponding unknown-neighbor count. Nonadjacent checks
are unaffected. This proves the identities inductively.

Initially the assertion about $x$ is its assumed agreement with $y$.
Suppose $x$ extends the current $K$, and the next step is $(j,i,b)$.
Every neighbor of $j$ except $i$ is already known. Since the $j$-th
parity of $x$ is zero,
\[
 x_i=\sum_{h\in\Nb(j)\setminus\{i\}}x_h
     =\sum_{h\in\Nb(j)\cap\operatorname{dom}K}K(h)=s(j)=b.
\]
Thus $x$ extends the enlarged assignment too. Induction proves the
claim for the entire trace, including inputs that have no completion
(for which the assertion is vacuous).
\end{proof}

A nonempty $S\subseteq\Var$ is a \emph{stopping set} when every check
has either zero or at least two neighbors in $S$ \cite{DiEtAl2002Stopping}.
\begin{prop}[Stopping sets]\label{prop:core-stopping}
A maximal peeling run leaves a nonempty unknown set precisely when that
remaining set is a stopping set. Peeling removes all initial erasures
if and only if the initial erasure set contains no stopping set.
Here ``removes'' is a structural statement, not acceptance of a codeword.
\end{prop}
\begin{proof}
A step exists exactly when some check has one unknown neighbor.
Consequently a nonempty residual has no possible step exactly when it
is a stopping set. Now let $S$ be a stopping subset of the initial
unknown set. If peeling removed a member of $S$, consider the first
such removal. At its witnessing check all members of $S$ are still
unknown; that check touches $S$ and hence has at least two unknown
neighbors, contradicting degree one. Thus $S$ survives every run.
Conversely, a nonempty residual of a maximal run is itself a stopping
subset of the initial erasure set, by the first part. These implications
establish the characterization independently of the choice of eligible
checks. A degree-zero row with syndrome one may still rule out any
completion; it does not change the structural argument.
\end{proof}

\section{Residual algebra and certificate checking}\label{sec:certificates}
For the final unknown set $U$ of a valid trace, order its coordinates
in increasing order, let $H_U$ be the corresponding column restriction
of $H$, and put $r=(s(j))_{j\in\Chk}$. A residual assignment
$\rho\in\Ftwo^{|U|}$ fills those coordinates, and $K\cup\rho$ denotes
the resulting full word.

\begin{lem}[Residual equivalence]\label{lem:core-residual}
For a valid trace from $y$, the map $\rho\mapsto K\cup\rho$ is a
bijection between solutions of $H_U\rho=r$ and codewords agreeing with
all visible coordinates of $y$.
\end{lem}
\begin{proof}
For a word extending $K$, its $j$-th parity is
$s(j)+\sum_{i\in\Nb(j)\cap U}x_i$, by
Lemma~\ref{lem:core-invariant}. This parity is zero exactly when the
$j$-th equation of $H_Ux_U=r$ holds. Hence every residual solution
produces a compatible codeword. In the other direction,
Lemma~\ref{lem:core-invariant} says that every compatible codeword
extends $K$; its restriction to $U$ therefore solves the residual
system. Restriction and filling are inverse functions.
\end{proof}

\begin{defi}[Certificate formats and independent checker]\label{def:certificates}
A success certificate consists of the tag $\mathsf{success}$, a trace
$T$, and a binary witness $\rho$ of length $|U|$ after replay. A failure
certificate consists of the tag $\mathsf{failure}$, a trace $T$, and a
witness $\omega\in\Ftwo^m$.
The checker first rejects a trace longer than $n$ (or stops a streaming
trace after its $(n+1)$-st entry). It then validates the tag, triple
structure, integer index bounds, binary values and witness lengths. It initializes the coherent
state from $(H,y)$ and replays every triple according to
Definition~\ref{def:peeling}; it rejects at the first invalid step.
For success it accepts exactly when $H_U\rho=r$, returning $K\cup\rho$.
For failure it accepts exactly when
\begin{equation}\label{eq:separator}
 \omega^\top H_U=0,\qquad \omega^\top r=1.
\end{equation}
The checker uses only the supplied instance and certificate; it does
not invoke the decoder. All rows are checked, including zero-degree
rows. A schema-valid partial trace is allowed in both branches.
\end{defi}

A hybrid producer runs maximal peeling, constructs the residual, and
performs exact Gaussian elimination over $\Ftwo$. A solution produces
a success certificate; an inconsistent row produces a failure
certificate as described below. When $U$ is empty, the residual still
has $m$ equations: success requires $r=0$. Removing all erasures does
not justify skipping this check.

\begin{thm}[Success soundness]\label{thm:core-success}
If the checker accepts a success certificate and returns $x$, then
$x\in\Ftwo^n$, $Hx=0$, and $x_i=y_i$ at every nonerased coordinate.
\end{thm}
\begin{proof}
The validated binary witness and assignment $K$ define a binary word.
Replay only assigns initially unknown coordinates, so $K$ extends
$K_0$ and channel agreement holds. The accepted residual equation and
Lemma~\ref{lem:core-residual} establish $Hx=0$.
\end{proof}

\begin{lem}[Finite-field alternative]\label{lem:core-alternative}
For $A\in\Ftwo^{m\times q}$ and $r\in\Ftwo^m$, exactly one of the
following exists: a solution $z$ of $Az=r$, or a vector $\omega$ with
$\omega^\top A=0$ and $\omega^\top r=1$.
\end{lem}
\begin{proof}
They cannot both exist, since then
$1=\omega^\top r=\omega^\top Az=0$.
Suppose there is no solution. Let $W$ be the span of the columns of
$A$. Then $r\notin W$. Choose a basis of $W$; appending $r$ gives a
linearly independent family, which extends to a basis of the finite
space $\Ftwo^m$. Define a linear functional $f$ on this basis to be
zero on the basis of $W$, one on $r$, and zero on all additionally
chosen basis vectors. Unique basis expansion defines $f$ on the
whole space and proves its linearity. If $e_j$ are the coordinate
basis vectors and $\omega_j=f(e_j)$, then for every vector $v$,
$f(v)=\sum_j v_j f(e_j)=\omega^\top v$. Therefore $f$ annihilates
every column of $A$, while $f(r)=1$, giving the required witness.
The argument also covers $q=0$. If $m=0$, the zero vector of $\Ftwo^q$ is always
a solution and the inconsistent case cannot arise.
\end{proof}

\begin{prop}[Constructing a failure witness]\label{prop:core-extraction}
Suppose elementary row operations transform $[A\mid r]$ to
$[R\mid c]=M[A\mid r]$, where $M\in\Ftwo^{m\times m}$ records the
same operations on the identity matrix. If $R$ has a zero row $h$
with $c_h=1$, then $\omega^\top=e_h^\top M$ satisfies the failure
conditions. Row reduction of an inconsistent system always finds
such a row.
\end{prop}
\begin{proof}
Taking row $h$ of the displayed equality gives
$(e_h^\top M)A=R_{h,*}=0$ and $(e_h^\top M)r=c_h=1$.
Row swaps and adding one row to another are invertible over $\Ftwo$,
so they preserve solvability. In echelon form, if every zero coefficient
row has right-hand side zero, set all free variables to zero and solve
the pivot equations from bottom to top. This constructs a solution:
each pivot is one, later pivot variables are already assigned, and
zero rows impose no further condition. Consequently an inconsistent
echelon system must have a zero coefficient row with right-hand side
one. Redundant rows and zero-column matrices cause no exception.
\end{proof}

\begin{thm}[Failure soundness and certificate completeness]\label{thm:core-failure}
An accepted failure certificate excludes every codeword agreeing with
$y$. Conversely, for every $(H,y)$ and every valid finite trace $T$,
there is an accepted certificate of exactly one status: success if a
compatible codeword exists, and failure otherwise.
\end{thm}
\begin{proof}
For soundness, acceptance supplies \eqref{eq:separator}. By
Lemma~\ref{lem:core-residual}, any compatible codeword would supply
a residual solution. This contradicts Lemma~\ref{lem:core-alternative}.
For completeness, if a compatible codeword exists, its restriction
to $U$ supplies a success witness by the residual bijection.
Otherwise the residual is inconsistent and the finite-field
alternative supplies a failure witness. Both witnesses pass the
specified checks, and the two statuses cannot both be accepted by
the same alternative. This is an existence statement for a suitable
witness, not a claim that every alleged witness is accepted.
\end{proof}

\begin{prop}[Multiplicity and exact recovery]\label{prop:core-nullity}
If $H_U\rho=r$ is consistent and has rank $a$, it has exactly
$2^{|U|-a}$ solutions. Hence the compatible completion is unique if
and only if $|U|-a=0$. Under a clean erasure channel the transmitted
word is one of these completions; acceptance alone does not identify
it when more than one completion exists.
\end{prop}
\begin{proof}
For one solution $\rho_0$, every solution is uniquely of the form
$\rho_0+v$ with $v\in\ker H_U$, because
$H_U(\rho+\rho_0)=r+r=0$. Rank--nullity gives dimension $|U|-a$ for
this kernel. A basis of a $d$-dimensional binary vector space gives a
bijection with $\Ftwo^d$, which has $2^d$ elements. Apply the residual
bijection. On the clean channel the transmitted codeword agrees
with the visible symbols, proving the last assertion.
\end{proof}

\section{Likelihood, examples and trust boundary}\label{sec:interpretation}
For erasure probability $\epsilon\in[0,1]$, define the memoryless BEC
transition probability by
\[
 W_\epsilon(y\mid x)=\prod_{i=1}^n w_\epsilon(y_i\mid x_i),\qquad
 w_\epsilon(a\mid b)=
 \begin{cases}
  \epsilon,&a=\starE,\\
  1-\epsilon,&a=b,\\
  0,&a\in\Ftwo\text{ and }a\ne b.
 \end{cases}
\]
Maximum likelihood means maximizing this quantity over $x\in C$;
it does not impose a prior distribution on codewords.
\begin{prop}[BEC maximum likelihood]\label{prop:core-ml}
For every $\epsilon\in[0,1]$, a compatible codeword maximizes
$W_\epsilon(y\mid x)$ over $C$. If $0<\epsilon<1$ and a compatible
codeword exists, the maximizers are exactly the compatible codewords.
\end{prop}
\begin{proof}
Let $u$ and $k=n-u$ be the numbers of erased and visible coordinates.
Every compatible word has likelihood
$c=\epsilon^u(1-\epsilon)^k\geq0$; every incompatible word has a
zero factor and therefore likelihood zero. Thus no candidate has
likelihood greater than a compatible word. For interior $\epsilon$,
$c>0$, so the incompatible candidates are not maximizers. The product
definition handles the endpoints, with an empty product equal to one.
An observation impossible at an endpoint may give all candidates
likelihood zero; the first assertion still holds, but the second
need not. No posterior conditional on such a zero-probability event
is asserted.
\end{proof}

\begin{exa}[Peeling, stopping, ambiguity and impossibility]\label{exa:core}
All indices in these calculations are one-based; the software uses
zero-based indices.
\begin{enumerate}[label=(\alph*),leftmargin=*]
\item For
$H=\left(\begin{smallmatrix}1&1&0\\0&1&1\end{smallmatrix}\right)$ and
$y=(0,\starE,\starE)$, initially $s=(0,0)$ and $d=(1,2)$.
The step $(1,2,0)$ changes these to $s=(0,0)$, $d=(0,1)$.
The step $(2,3,0)$ leaves $U=\varnothing$ and $r=(0,0)$,
so the empty residual assignment certifies $x=(0,0,0)$.
\item For $H=(1\ 1)$ and $y=(\starE,\starE)$, the unknown set is
a stopping set and the empty trace is maximal. The residual equation
$\rho_1+\rho_2=0$ has solutions $(0,0)$ and $(1,1)$, rank one and
nullity one. Each is an ML completion, and neither is selected by
likelihood alone.
\item For
$H=\left(\begin{smallmatrix}1&1&0\\1&0&1\end{smallmatrix}\right)$,
$y=(\starE,0,1)$ and $T=\varnothing$, the residual is
$A=(1,1)^\top$, $r=(0,1)^\top$.
The witness $\omega=(1,1)^\top$ gives $\omega^\top A=0$ and
$\omega^\top r=1$. It certifies impossibility before any peeling.
Alternatively resolving $x_1=0$ from row one leaves no residual
variables and right-hand side $(0,1)$; $\omega=(0,1)^\top$ then
certifies the remaining degree-zero contradiction.
\item For $H=(1\ 1)$, a transmission of $(0,0)$ followed by one
visible flip and one erasure can yield $y=(1,\starE)$.
Peeling validly returns $(1,1)$, which passes both parity and channel
checks. The certificate is valid, but the transmitted word was not
recovered. Corruption need not make the observed constraints
inconsistent.
\end{enumerate}
\end{exa}

The trusted instance is the supplied pair $(H,y)$, including its
validated graph representation. The mathematical checker assumes
exact binary arithmetic; an executable deployment also trusts its
parser, checker implementation, interpreter/compiler and runtime.
The decoder and its residual solver need not be trusted. A Lean
proof of a mathematical function does not by itself verify an
independently written Python program or its execution environment.

An accepted trace certifies that the declared sequence is permitted
by the specification. It does not attest that an untrusted producer
actually followed this sequence internally. A valid partial trace,
or a permutation of independent steps that remains valid, must be
accepted. Malformed certificates, incompatible words and invalid
steps must be rejected. Rejection of a malformed failure witness
alone does not establish that a completion exists.

If the supplied matrix is corrupted, acceptance proves only the
constraints of that supplied matrix, not those of an intended
external code. Authenticating the matrix or channel observations
is outside this model. On inputs containing visible flips,
compatibility is an algebraic property; the BEC likelihood theorem
is not a guarantee about the physical channel that generated them.

\section{Representation, verification cost and certificate size}\label{sec:cost}
Assume a validated graph stored as check-to-variable and
variable-to-check adjacency lists with no duplicate incidences.
Use a word-RAM with words large enough for indices and counters,
constant-time binary arithmetic and array access. Graph validation
and construction from a dense input matrix require $\Theta(mn)$
entry inspections; they are separate from the following per-input
bounds. Storing a matrix densely does not become a sparse algorithm
merely by calling a pivot heuristic ``sparse.''

\begin{prop}[Sparse checking upper bound]\label{prop:core-cost}
A success or failure certificate can be checked in
$O(n+m+|\Edges|)$ word operations and $O(n+m+|\Edges|)$ space,
including graph storage. Direct verification of a binary full word
against $H$ and $y$ has the same asymptotic upper bound and requires
no trace.
\end{prop}
\begin{proof}
Reading $y$ and initializing variable status costs $O(n)$; initializing
check arrays costs $O(m)$. Visiting each incidence once computes
syndromes, residual degrees and the XOR of the indices of each
check's unknown neighbors, at cost $O(|\Edges|)$.
When a residual degree is one, this last XOR identifies its unique
unknown index in constant time, even if that index is zero.
The checker validates each triple using these arrays, then visits
the adjacency list of its newly assigned variable to update the
three check arrays. No variable is assigned twice, so the total
number of such incidence visits is at most $|\Edges|$, and there are
at most $n$ valid triples. An excessive trace is rejected after at
most $n+1$ inspected entries; malformed dimensions are rejected
before witness processing.

For success, fill the remaining coordinates from the witness and
accumulate their contributions across the unknown incidences,
checking all $m$ residual equations. For failure, allocate one binary
accumulator per unknown variable, and visit each residual incidence
$(j,i)$ to add $\omega_j$ to its accumulator; separately compute
$\sum_j\omega_jr_j$. These procedures cost $O(n+m+|\Edges|)$.
The arrays have $O(n+m)$ entries and the graph has $O(|\Edges|)$
incidences. Direct checking reads the candidate and $y$ and computes
each parity by a single pass over the check adjacency lists.
\end{proof}

The bound simplifies to $O(n+|\Edges|)$ only under an assumption such
as $m=O(n+|\Edges|)$. It is an upper bound for this representation,
not an optimality theorem for every possible certificate format or
preprocessed matrix. The bit cost of reading indices is not constant
in a bit-complexity model. A parser that first materializes an
arbitrarily long hostile input additionally pays for those bytes;
the graph bound assumes a parsed certificate of bounded format or a
streaming parser that enforces these limits.

With $t=|T|$, an index-encoded trace occupies
$O(t(1+\log(m+1)+\log(n+1)))$ bits, plus a length/tag description.
A success witness adds $|U|$ bits in canonical variable order; a
failure witness adds $m$ bits. For $m=O(n)$ and $t\leq n$, these are
$O(n\log(n+1))$ and $O(n\log(n+1)+m)$ respectively.
A direct success witness is simply an $n$-bit word. It proves parity
and visible-coordinate agreement without proving a declared peeling
trace. Consequently trace replay purchases an additional explanatory
property, not a stronger final-word validity guarantee. For an
empty trace, a failure witness for the initial erasure system already
certifies impossibility without any peeling history.

Certificate \emph{construction} has a different cost. A straightforward
dense elimination of an $m$ by $q$ residual, with $m$ extra witness
columns, takes $O(q+m(q+m)(1+\min(m,q)))$ elementary bit operations,
including initialization, row-array updates and pivot searches; its storage is
$O(q+m(q+m))$ bits apart from indexing, including the returned solution. Packed machine-word rows can
improve constants and word-operation bounds. Peeling may reduce $q$,
but no universal speedup over GE-only follows, and recording the
row transformations can itself be expensive. Our comparisons use
the same algebraic engine to keep that distinction observable.
